% GENERATED FILE. Edit canonical lessons, then run npm run generate:books.
% canonical-source-hash: fnv1a64:6fd670e47327a204
% canonical-lessons: HI-C421-basti, HI-C421-qasba, HI-C421-basinda, HI-C421-paros, HI-C421-dehat, HI-R421-review-414-421

\chapter{Settlements and Towns}
\label{ch:hi-settlements-and-towns}
\hlchaptermodality{421}

\begin{chapteropening}
\textbf{By the end of this chapter:} \emph{I can talk about a settlement, a small town and the countryside, and the people who live there.}

Complete the last lesson of chapter 421: I can talk about a settlement, a small town and the countryside, and the people who live there.
\end{chapteropening}

\section[\texorpdfstring{\hi{बस्ती} (bastī)}{बस्ती (bastī)}]{\hi{बस्ती} (bastī) — a settlement; a poor quarter}

\label{lesson:HI-C421-basti}



\begin{quote}
Before the new one: say the Hindi for to roll, then the Hindi for to dig into.
\end{quote}

\subsection*{You'll want to know: \hi{बस्ती}}
\textbf{\hi{बस्ती}} — \emph{bastī} — \textquotedblleft{}a settlement; a poor quarter\textquotedblright{}.

\begin{quote}
\hi{नदी} \hi{के} \hi{पार} \hi{एक} \hi{छोटी} \hi{बस्ती} \hi{है।}
\end{quote}

\begin{quote}
\emph{nadī ke pār ek choṭī bastī hai.}
\end{quote}

\begin{quote}
There is a small settlement across the river.
\end{quote}

Feminine.

\begin{cousinweb}
\textbf{\hi{बस्ती}} (\emph{bastī}) is usually traced to Sanskrit \emph{vasati}, \textquotedblleft{}staying, a dwelling-place\textquotedblright{} (Monier-Williams), from the root \emph{vas}, \textquotedblleft{}to dwell\textquotedblright{}. A \hi{बस्ती} can be a hamlet or a crowded poor quarter of a city.
\end{cousinweb}

\begin{grammarlens}[title={What you've built}]
One more word for getting around.
\end{grammarlens}

\subsection*{Your turn}
\begin{itemize}
\raggedright
  \item \emph{Say it:} \emph{bastī}
  \item \emph{Say it:} \emph{bastī}, once more
  \item \emph{Say it:} \emph{nadī ke pār ek choṭī bastī hai.} — There is a small settlement across the river.
  \item \emph{Recall:} say the Hindi for to roll, then the Hindi for to dig into, then say \emph{bastī} again
\end{itemize}

\begin{tcolorbox}[breakable,colback=teal!4,colframe=teal!35!black,title={Before you move on}]
What does \hi{बस्ती} mean? (\textquotedblleft{}A settlement; a poor quarter\textquotedblright{}.) Say it once more.
\end{tcolorbox}
\section[\texorpdfstring{\hi{क़स्बा} (qasbā)}{क़स्बा (qasbā)}]{\hi{क़स्बा} (qasbā) — a small town}

\label{lesson:HI-C421-qasba}



\begin{quote}
Before the new one: say the Hindi for to dig into, then the Hindi for a settlement.
\end{quote}

\subsection*{You'll want to know: \hi{क़स्बा}}
\textbf{\hi{क़स्बा}} — \emph{qasbā} — \textquotedblleft{}a small town\textquotedblright{}.

\begin{quote}
\hi{मेरे} \hi{नाना} \hi{एक} \hi{छोटे} \hi{क़स्बे} \hi{में} \hi{रहते} \hi{हैं।}
\end{quote}

\begin{quote}
\emph{mere nānā ek choṭe qasbe meṁ rahte haiṁ.}
\end{quote}

\begin{quote}
My grandfather lives in a small town.
\end{quote}

Masculine.

\begin{cousinweb}
\textbf{\hi{क़स्बा}} (\emph{qasbā}) is usually traced to Arabic \emph{qaṣaba}, \textquotedblleft{}a chief town, a fortress\textquotedblright{}, which came through Persian; in Hindi a \hi{क़स्बा} is bigger than a village and smaller than a city.
\end{cousinweb}

\begin{grammarlens}[title={What you've built}]
One more word for getting around.
\end{grammarlens}

\subsection*{Your turn}
\begin{itemize}
\raggedright
  \item \emph{Say it:} \emph{qasbā}
  \item \emph{Say it:} \emph{qasbā}, once more
  \item \emph{Say it:} \emph{mere nānā ek choṭe qasbe meṁ rahte haiṁ.} — My grandfather lives in a small town.
  \item \emph{Recall:} say the Hindi for to dig into, then the Hindi for a settlement, then say \emph{qasbā} again
\end{itemize}

\begin{tcolorbox}[breakable,colback=teal!4,colframe=teal!35!black,title={Before you move on}]
What does \hi{क़स्बा} mean? (\textquotedblleft{}A small town\textquotedblright{}.) Say it once more.
\end{tcolorbox}
\section[\texorpdfstring{\hi{बाशिंदा} (bāśindā)}{बाशिंदा (bāśindā)}]{\hi{बाशिंदा} (bāśindā) — a resident, an inhabitant}

\label{lesson:HI-C421-basinda}



\begin{quote}
Before the new one: say the Hindi for a settlement, then the Hindi for a small town.
\end{quote}

\subsection*{You'll want to know: \hi{बाशिंदा}}
\textbf{\hi{बाशिंदा}} — \emph{bāśindā} — \textquotedblleft{}a resident, an inhabitant\textquotedblright{}.

\begin{quote}
\hi{मैं} \hi{इस} \hi{शहर} \hi{का} \hi{पुराना} \hi{बाशिंदा} \hi{हूँ।}
\end{quote}

\begin{quote}
\emph{maiṁ is śahar kā purānā bāśindā hūṁ.}
\end{quote}

\begin{quote}
I am a long-time resident of this city.
\end{quote}

Masculine.

\begin{cousinweb}
\textbf{\hi{बाशिंदा}} (\emph{bāśindā}) is usually traced to Persian \emph{bāšinda}, \textquotedblleft{}one who stays, an inhabitant\textquotedblright{}.
\end{cousinweb}

\begin{grammarlens}[title={What you've built}]
One more word for getting around.
\end{grammarlens}

\subsection*{Your turn}
\begin{itemize}
\raggedright
  \item \emph{Say it:} \emph{bāśindā}
  \item \emph{Say it:} \emph{bāśindā}, once more
  \item \emph{Say it:} \emph{maiṁ is śahar kā purānā bāśindā hūṁ.} — I am a long-time resident of this city.
  \item \emph{Recall:} say the Hindi for a settlement, then the Hindi for a small town, then say \emph{bāśindā} again
\end{itemize}

\begin{tcolorbox}[breakable,colback=teal!4,colframe=teal!35!black,title={Before you move on}]
What does \hi{बाशिंदा} mean? (\textquotedblleft{}A resident, an inhabitant\textquotedblright{}.) Say it once more.
\end{tcolorbox}
\section[\texorpdfstring{\hi{पड़ोस} (paṛos)}{पड़ोस (paṛos)}]{\hi{पड़ोस} (paṛos) — the neighbourhood; next door}

\label{lesson:HI-C421-paros}



\begin{quote}
Before the new one: say the Hindi for a small town, then the Hindi for a resident.
\end{quote}

\subsection*{You'll want to know: \hi{पड़ोस}}
\textbf{\hi{पड़ोस}} — \emph{paṛos} — \textquotedblleft{}the neighbourhood; next door\textquotedblright{}.

\begin{quote}
\hi{हमारे} \hi{पड़ोस} \hi{में} \hi{एक} \hi{नई} \hi{दुकान} \hi{खुली} \hi{है।}
\end{quote}

\begin{quote}
\emph{hamāre paṛos meṁ ek naī dukān khulī hai.}
\end{quote}

\begin{quote}
A new shop has opened in our neighbourhood.
\end{quote}

Masculine.

\begin{cousinweb}
\textbf{\hi{पड़ोस}} (\emph{paṛos}) is a tadbhava of Sanskrit \emph{prativeśa}, \textquotedblleft{}neighbouring; a neighbour\textquotedblright{} (CDIAL 8598), which lists Hindi \emph{paṛos}, \textquotedblleft{}neighbourhood\textquotedblright{}. \hi{पड़ोसी}, a neighbour, is built on it.
\end{cousinweb}

\begin{grammarlens}[title={What you've built}]
One more word for getting around.
\end{grammarlens}

\subsection*{Your turn}
\begin{itemize}
\raggedright
  \item \emph{Say it:} \emph{paṛos}
  \item \emph{Say it:} \emph{paṛos}, once more
  \item \emph{Say it:} \emph{hamāre paṛos meṁ ek naī dukān khulī hai.} — A new shop has opened in our neighbourhood.
  \item \emph{Recall:} say the Hindi for a small town, then the Hindi for a resident, then say \emph{paṛos} again
\end{itemize}

\begin{tcolorbox}[breakable,colback=teal!4,colframe=teal!35!black,title={Before you move on}]
What does \hi{पड़ोस} mean? (\textquotedblleft{}The neighbourhood; next door\textquotedblright{}.) Say it once more.
\end{tcolorbox}
\section[\texorpdfstring{\hi{देहात} (dehāt)}{देहात (dehāt)}]{\hi{देहात} (dehāt) — the countryside; villages}

\label{lesson:HI-C421-dehat}



\begin{quote}
Before the new one: say the Hindi for a resident, then the Hindi for next door.
\end{quote}

\subsection*{You'll want to know: \hi{देहात}}
\textbf{\hi{देहात}} — \emph{dehāt} — \textquotedblleft{}the countryside; villages\textquotedblright{}.

\begin{quote}
\hi{मेरे} \hi{दादा} \hi{देहात} \hi{में} \hi{रहते} \hi{हैं।}
\end{quote}

\begin{quote}
\emph{mere dādā dehāt meṁ rahte haiṁ.}
\end{quote}

\begin{quote}
My grandfather lives in the countryside.
\end{quote}

Masculine.

\begin{cousinweb}
\textbf{\hi{देहात}} (\emph{dehāt}) is usually traced to Persian \emph{dehāt}, \textquotedblleft{}villages\textquotedblright{}, the plural of \emph{deh}, \textquotedblleft{}a village\textquotedblright{}; in Hindi it is the countryside as a whole, set against the town.
\end{cousinweb}

\begin{grammarlens}[title={What you've built}]
One more word for getting around.
\end{grammarlens}

\subsection*{Your turn}
\begin{itemize}
\raggedright
  \item \emph{Say it:} \emph{dehāt}
  \item \emph{Say it:} \emph{dehāt}, once more
  \item \emph{Say it:} \emph{mere dādā dehāt meṁ rahte haiṁ.} — My grandfather lives in the countryside.
  \item \emph{Recall:} say the Hindi for a resident, then the Hindi for next door, then say \emph{dehāt} again
\end{itemize}

\begin{tcolorbox}[breakable,colback=teal!4,colframe=teal!35!black,title={Before you move on}]
What does \hi{देहात} mean? (\textquotedblleft{}The countryside; villages\textquotedblright{}.) Say it once more.
\end{tcolorbox}
\section[(dialogue)]{Looking back: from Bathrooms and Comfort to Settlements and Towns}

\label{lesson:HI-R421-review-414-421}



\begin{quote}
Say the Hindi for a bathroom and a toilet.
\end{quote}

\subsection*{Your turn}
Say these aloud:

\begin{itemize}
\raggedright
  \item a bathroom, a toilet, a water heater, comfortable, to wake (someone)
  \item a manager, a register, a booking, a place to stay, to send for
  \item a tent (for camping), to put to sleep, a roost, to spend (time), to pass (time)
  \item a direction, a side (of a place), towards, a route, the way something faces
  \item close by, near (formal), across, beside, to deliver
  \item a T-junction, a round, a mark left behind, a building, a signboard
  \item a slope, a corridor, to push in, to roll, to dig into
  \item a settlement, a small town, a resident, next door, the countryside
\end{itemize}

\begin{tcolorbox}[breakable,colback=teal!4,colframe=teal!35!black,title={Before you move on}]
A bathroom? (\textbf{\hi{गुसलख़ाना}}, \emph{gusalkhānā}.) A toilet? (\textbf{\hi{शौचालय}}, \emph{śaucālay}.) A water heater? (\textbf{\hi{गीज़र}}, \emph{gīzar}.) Comfortable? (\textbf{\hi{आरामदेह}}, \emph{ārāmdeh}.) To wake (someone)? (\textbf{\hi{जगाना}}, \emph{jagānā}.) A manager? (\textbf{\hi{मैनेजर}}, \emph{mainejar}.) A register? (\textbf{\hi{रजिस्टर}}, \emph{rajisṭar}.) A booking? (\textbf{\hi{बुकिंग}}, \emph{buking}.) A place to stay? (\textbf{\hi{ठिकाना}}, \emph{ṭhikānā}.) To send for? (\textbf{\hi{मँगाना}}, \emph{mãgānā}.) A tent (for camping)? (\textbf{\hi{ख़ेमा}}, \emph{khemā}.) To put to sleep? (\textbf{\hi{सुलाना}}, \emph{sulānā}.) A roost? (\textbf{\hi{बसेरा}}, \emph{baserā}.) To spend (time)? (\textbf{\hi{गुज़ारना}}, \emph{guzārnā}.) To pass (time)? (\textbf{\hi{बिताना}}, \emph{bitānā}.) A direction? (\textbf{\hi{दिशा}}, \emph{diśā}.) A side (of a place)? (\textbf{\hi{तरफ़}}, \emph{taraf}.) Towards? (\textbf{\hi{ओर}}, \emph{or}.) A route? (\textbf{\hi{मार्ग}}, \emph{mārg}.) The way something faces? (\textbf{\hi{रुख़}}, \emph{rukh}.) Close by? (\textbf{\hi{नज़दीक}}, \emph{nazdīk}.) Near (formal)? (\textbf{\hi{निकट}}, \emph{nikaṭ}.) Across? (\textbf{\hi{पार}}, \emph{pār}.) Beside? (\textbf{\hi{बग़ल}}, \emph{baġal}.) To deliver? (\textbf{\hi{पहुँचाना}}, \emph{pahũcānā}.) A T-junction? (\textbf{\hi{तिराहा}}, \emph{tirāhā}.) A round? (\textbf{\hi{चक्कर}}, \emph{cakkar}.) A mark left behind? (\textbf{\hi{निशान}}, \emph{niśān}.) A building? (\textbf{\hi{इमारत}}, \emph{imārat}.) A signboard? (\textbf{\hi{बोर्ड}}, \emph{borḍ}.) A slope? (\textbf{\hi{ढलान}}, \emph{ḍhalān}.) A corridor? (\textbf{\hi{गलियारा}}, \emph{galiyārā}.) To push in? (\textbf{\hi{घुसना}}, \emph{ghusnā}.) To roll? (\textbf{\hi{लुढ़कना}}, \emph{luṛhaknā}.) To dig into? (\textbf{\hi{गड़ना}}, \emph{gaṛnā}.) A settlement? (\textbf{\hi{बस्ती}}, \emph{bastī}.) A small town? (\textbf{\hi{क़स्बा}}, \emph{qasbā}.) A resident? (\textbf{\hi{बाशिंदा}}, \emph{bāśindā}.) Next door? (\textbf{\hi{पड़ोस}}, \emph{paṛos}.) The countryside? (\textbf{\hi{देहात}}, \emph{dehāt}.)
\end{tcolorbox}
